\documentclass[twocolumn]{article}
 
\usepackage{epstopdf}
\usepackage[caption=false]{subfig}
\usepackage[numbers,sort&compress]{natbib}
\bibpunct[, ]{[}{]}{,}{n}{,}{,}
\renewcommand\bibfont{\fontsize{10}{12}\selectfont}

\usepackage{amsmath,amssymb,amsfonts}
\usepackage{graphicx}
\usepackage{booktabs}
\usepackage{xcolor}
\usepackage{ragged2e}
\setcounter{secnumdepth}{0}
\setlength{\parindent}{0pt}
\AtBeginDocument{
\setlength{\abovedisplayskip}{5pt}
\setlength{\belowdisplayskip}{5pt}
\setlength{\abovedisplayshortskip}{3pt}
\setlength{\belowdisplayshortskip}{3pt}
}

\begin{document}

\twocolumn[{
\vspace{-2cm} 
\raggedright

{\Large\bfseries
Deep Learning-Based Prediction of Nifty BeES Price Movement Using Financial News from Indian Media\par}

\vspace{6pt}

Vishnupriya Thimmarayudu\textsuperscript{a}, 
Advaith Nanda Kishore\textsuperscript{a}, 
Busi Reddy Nischitha Reddy\textsuperscript{a}, 
Aman Joy Tirkey\textsuperscript{a}, 
Shobharani D.~A.~\textsuperscript{a}\par

\vspace{4pt}

\textsuperscript{a}GITAM University, Bengaluru, Karnataka, India \\
Email: vthimma@gitam.in\par

\vspace{6pt}

\justifying

\noindent\textbf{Abstract:}
The following paper aims to present a AI based stock trading system designed for the Indian Stock markets (NSE and BSE). The system built utilizes multiple machine learning methodologies, mainly Long Short Term Memory or LSTM for predicting price, Random Forest Algorithm for assessing the risks, VADER and TextBlob for sentiment analysis, Retrieval Augmented Generation chatbot for recommendation with justifications. The approach used in building the system combines multiple techniques where in technical indicators, sentiment analysis from news and social media, real time market data for best possible insights for intraday as well as swing trading. Random Forest achieves 90--94\% accuracy (CV: 0.9429). Sentiment analysis runs on real time data, processing the news from multiple sources, keeping a consistent score. LSTM predicts next-day price direction. This complete system is deployed as a full stack web application termed as TradeLens, having 51 trained models, capable of supporting not only the Indian market but also international markets.

\vspace{6pt}

\noindent\textbf{Keywords:}
Stock Market Prediction; Long Short-Term Memory (LSTM); Random Forest; Sentiment Analysis; Retrieval-Augmented Generation (RAG); Financial AI; Explainable AI (XAI)

\vspace{10pt}
}]


%=======================
\section{Introduction}
%=======================
The Indian financial market is extremely dynamic and dense in terms of information. It consists of the National Stock Exchange and Bombay Stock Exchange. With rise in ease of access for investors through mobile brokerage applications, the number of participants has increased significantly. However, being able to interpret the market movements still remains a challenge due to the high level of complexity of financial data and limited tools available for the investors. Retail traders feel overwhelmed due to excessive volumes of data which turns into noise. Rapid changes in news, social media discussions, and price data leads to reduced accuracy in decision making.

Traditional market prediction models, generally depend on historical price patterns and technical indicators which are not sufficient alone for being able to capture sentiment driven movements in market caused due to numerous reasons such as policy changes, macroeconomic changes or financial news \cite{bollen2011twitter}. Whereas, purely depending on qualitative approach which is based on financial news can affect the quantitative aspects negatively while taking trading decisions. TradeLens was designed to ensure these challenges are mitigated with the help of an integrated artificial intelligence framework which consists of deep learning based forecasting of price, better risk assessment and sentiment analysis. It leverages Retrieval Augmented Generation(RAG) architecture \cite{lewis2020rag}, hence providing explanations which can be easily interpreted by investors and understand the reasoning to every trading signal provided. 

With the improvements in the current deep learning domain, it has allowed models to work better with financial time series data. LSTM network, a type of recurrent neural network, helps to mitigate vanishing gradient problem \cite{hochreiter1997lstm}, making it very useful to spot temporal dependencies in long range. The proposed model, TradeLens uses LSTM with a lookback window of 60 days of trading, ensuring the prediction is made keeping both technical indicators and historical price data in consideration. 

To improve the strength of the model in volatile market conditions, the system consists of an ensemble based risk classification module using Random Forest(RF) \cite{breiman2001rf}. Random Forest creates several decision trees during training and finalizes the prediction decision through voting.

\subsection{Motivation}

Retail investors have shown high participation over time which brings the need for an AI based analytical tool that is easily available. SEBI focuses on transparency, and it is noticeable that many models till date function as black box systems, which lowers confidence of the investors before taking decisions. Hence, to address this issue, TradeLens being a hybrid system is a combination of LSTM, Random Forest, sentiment analysis, and Retrieval Augmented Generation (RAG) hence providing better interpretability to investors along with reliability in predictions, hence boosting investors confidence in decision making.

\subsection{Contributions}
\begin{enumerate}
    \item A hybrid AI driven framework which performs faster with both LSTM based temporal forecasting and Random forest based risk classification, with a high risk assessment accuracy of 94.3\%.
    \item A highly precise sentiment pipeline using VADER and TextBlob architecture to ensure it meets the high complexity of market news signals and social media signals.
    \item A professional level RAG-based AI assistant that is integrated to provide justifications to each signal provided that are easy to understand with factual recommendations.
\end{enumerate}

%========================
\section{Related Work}
%========================

\subsection{Deep Learning Architectures for Financial Time-Series}

The prediction of stock market has moved from traditional models like ARIMA and GARCH to powerful deep learning models which can capture non linear dynamic changes. Long Short Term Memory (LSTM) networks~\cite{hochreiter1997lstm} prevent vanishing gradient problems and also benefit with long term temporal dependencies. Previous studies have shown that LSTM can perform better than Random Forest~\cite{breiman2001rf} and neural networks in directional prediction.

Studies recently have focused on markets where volatility is very noticeable. Works by Kallimath et al.~\cite{kallimath2024nifty} have confirmed how effective LSTM based models are for forecasting price for Nifty. Using Bi-LSTM and hybrid systems have shown better results as proven in recent comparative studies~\cite{barua2024comparative}, proving the suitability of recurrent models for financial market predictions. Paul and Das~\cite{paul2024bilstm} further confirmed that BiLSTM consistently outperforms standard RNN architectures on Indian market indices.

\subsection{Sentiment Analysis in Financial Markets}

Sentiment analysis is a crucial component in financial systems due to the weight it carries in driving prices in the market. Bollen et al.~\cite{bollen2011twitter} have proven how social media affects markets. VADER, introduced by Hutto and Gilbert~\cite{hutto2014vader} shows how efficiency in sentiment scoring can be performed for social media texts, however it lacked contextual understanding.

Recent studies showed that integration of sentiment with machine learning improves accuracy in predicting markets. Surveys recently show a shift towards integration of sentiment signals along with quantitative models~\cite{chandra2025survey,bhardwaj2015india}, whereas, Garg et al.~\cite{garg2024ipo} focus on predictive value of sentiments.

\subsection{Hybrid and Ensemble Models}

Hybrid and ensemble approaches have increased as they have the ability to balance the variance with bias. Zhu and Wu~\cite{zhu2024hybrid} showed hybrid optimization for better feature selection. Previous work has addressed class imbalance in financial datasets. Ouf et al.~\cite{ouf2024sentiment} showed that LSTM combined with ensemble models and sentiment features can improve performance in volatile markets. 

Comparative studies show how hybrid models can outperform single model based approaches~\cite{maqbool2024sentiment}. Hence a regime filtering strategy, where Random Forest is a gating mechanism for prediction made by LSTM, improves the reliability in high volatility.

\subsection{Explainable AI in Finance}

Interpretability has always been a challenge in deep learning models with respect to financial domains. Bussmann et al.~\cite{bussmann2020explainability} show the importance interpretability plays for compliance with the regulatory. Retrieval-Augmented Generation (RAG), which was introduced by Lewis et al.~\cite{lewis2020rag} ensures solid explanations.

Chauhan et al.~\cite{chauhan2024augmented} integrated human insights with deep learning to ensure trust and transparency. TradeLens ensures RAG based explanation is provided for complete transparency on every signal of prediction provided.

%========================
\section{System Architecture and Data Pipeline}
%========================

\begin{figure}[!t]
\centering
\includegraphics[width=\columnwidth]{figures/FIG 10 .png}
\caption{System Architecture}
\label{fig:system_architecture}
\end{figure}


\subsection{System Architecture}

TradeLens integrates data ingestion, preprocessing, hybrid machine learning inference, frontend visualization layers as seen in Figure~\ref{fig:system_architecture}.

\begin{itemize}

\item \textbf{Data Ingestion Layer:}  
This layer focuses on collection of raw data from external sources. Datasets like market price, OHLCV values are sourced from Yahoo Finance~\cite{yahoofinance2024}. Finance news and social media signals are all obtained from APIs. These provide numerical as well as textual data which is required for predicting markets.

\item \textbf{Preprocessing Pipeline:}  
This layer converts the inputs into features that are structures and suitable for machine learning models. Panda dataframes are responsible to organize financial data, which then is followed by cleaning of data, normalization and feature transformation. RSI and MACD are some indicators which are computed and the results are scaled for stable training and inference of the model.

\item \textbf{Hybrid AI Engine:}  
Hybrid AI engine is the core of the entire architecture. LSTM networks analyze the price data sequentially to predict movements. Parallelly, Random Forest evaluates the volatility of market and risk levels. Combination of the above improves the reliability of signals.

\item \textbf{Backend Processing Layer:}  
Backend uses Flask REST API that ensures execution, API request, and scheduling of tasks occurs smoothly. Job scheduler updates based on market data, which ensures that no manual intervention is required during the operation of the system.

\item \textbf{Frontend Visualization Layer:}  
Frontend is React-based aiming for clean and easy to use interface for the users without the compromise of useful data for the users. It has built in Streamlit for monitoring the system status and managing the outputs in the model.

\end{itemize}


\subsection{Data Acquisition and Processing}

This pipeline is responsible for collecting financial and sentiment data from various sources. Datasets like market price, OHLCV values are sourced from Yahoo Finance~\cite{yahoofinance2024}. Finance news and social media signals are all obtained from APIs. \texttt{ThreadPoolExecutor} ensures the latency is reduced. Temporal alignment is maintained using the time based filtering to ensure synchronization of market and news data within the trading hours. A caching mechanism using JSON benefits reliability and reduces dependency solely on APIs.

After data collection, the data have to be preprocessed, which involves cleaning, normalization, and feature engineering. RSI and MACD are computed and scaled for stability in performance of the model.

The processed data are passed through a hybrid inference framework. LSTM predicts the direction by capturing the temporal dependencies and Random Forest checks on the market volatility and risk levels, both combining to give a powerful output.

For better interpretability, Retrieval Augmented Generation (RAG) module is utilized. Financial documents are embedded with the help of HuggingFace \textit{all-MiniLM-L6-v2} model and is stored in FAISS vector database. Gemini 1.5 Flash is used to generate human-readable explanations for better trading decisions.

\subsection{Dataset and Reproducibility Settings}

The model is trained and evaluated on Nifty BeES (NSE: NIFTYBEES) daily OHLCV data from January 2020 to December 2024, comprising approximately 1,200 trading days, sourced from Yahoo Finance~\cite{yahoofinance2024}. A strict temporal split is applied: 80\% for training (January 2020 to December 2023) and 20\% for testing (January 2024 to December 2024), ensuring no data leakage. Model performance is reported using 5-fold walk-forward time-series cross-validation on the training set. All experiments use random seed = 42, Python 3.10, TensorFlow 2.15, and Scikit-learn 1.4. All trained model weights are serialized using Joblib for full reproducibility~\cite{amruth2025stock}.

%========================
\section{Proposed Trading Inference Framework}
%========================

\subsection{Protocol Overview}

TradeLens consists of a multi stage hybrid protocol which provides reliable trading signals. It has three main components : LSTM based temporal analysis, Random Forest classifier which classifies the risk based on volatility and RAG based modules which ensure interpretability for users. 

\subsection{Hybrid Inference Phase}

Inference pipelines work sequentially, where the first step is, OHLCV data is obtained using the data acquisition engine, then, LSTM model predicts the direction of the price (Equation~\eqref{eq:lstm}), parallelly the Random Forest classifier ensures classification of risk (Equation~\eqref{eq:risk}). The sentiment features are also derived from news and other sources like social media which are eventually aggregated (Equation~\eqref{eq:sentiment}). The unified trading score (Equation~\eqref{eq:trading_score}) then determines the final trading signal (Equation~\eqref{eq:signal}). All combine to give better results.

%========================
\subsection{Mathematical Formulation of Hybrid Inference}
%========================

The proposed framework integrates LSTM-based prediction, sentiment fusion, and Random Forest-based risk classification into a unified decision function.

\begin{equation}
\hat{y}_{t+1} = \text{LSTM}(X_t)
\label{eq:lstm}
\end{equation}
where $\hat{y}_{t+1}$ represents the predicted price or directional movement obtained from the LSTM model using input features $X_t$.

\begin{equation}
S_t = \alpha \cdot \nu_t + (1-\alpha)\cdot b_t
\label{eq:sentiment}
\end{equation}
where $S_t$ denotes the aggregated sentiment score combining VADER ($\nu_t$) and TextBlob ($b_t$) polarity values.

\begin{equation}
R_{\text{risk}} = \text{RF}(X_t)
\label{eq:risk}
\end{equation}
where $R_{\text{risk}}$ represents the risk level estimated using the Random Forest classifier based on market features.

\begin{equation}
\mathcal{T}_t = w_1 \cdot \text{sign}(\hat{y}_{t+1}) + w_2 \cdot S_t - w_3 \cdot R_{\text{risk}}
\label{eq:trading_score}
\end{equation}
where $\mathcal{T}_t$ is the unified trading score integrating directional prediction, sentiment, and risk with weights $w_1, w_2,$ and $w_3$.

\begin{equation}
\text{Signal} =
\begin{cases}
\text{Buy}, & \mathcal{T}_t > \theta_b \\
\text{Sell}, & \mathcal{T}_t < \theta_s \\
\text{Hold}, & \text{otherwise}
\end{cases}
\label{eq:signal}
\end{equation}
where $\theta_b$ and $\theta_s$ are predefined thresholds used to generate the final trading decision.

\subsection{Feature Representation}

Predictions made are represented individually to the user to ensure clarity. The four components are LSTM price prediction, risk categorization, sentiment scores of news and social media sources, justification of the trading signal provided.

\subsection{Verification and Signal Generation}

The signal is finally validated using a series of constraints. Directional consistency is crosschecked using sign alignment between price changes and the sentiments at the given moment. Risk gating ensures that the signals that exceed a given volatility threshold are rejected. Temporal validity is ensured by restriction of outdated inputs. 

%========================
\section{Performance Evaluation}
%========================

\subsection{Evaluation Metrics}

The evaluation of the performance consists of metrics such as: Accuracy, Precision, Recall, and F1-score. They are obtained from the confusion matrix components: true positives (TP), true negatives (TN), false positives (FP), and false negatives (FN)~\cite{sokolova2020classification}.

\subsection{Quantitative Results}

Quantitative results of different models are as shown in Table~\ref{tab:model_results} using Accuracy, Precision, Recall, and F1-score.

Baseline models like Linear Regression and Random Forest are always outperformed by deep learning methods in every metrics. This shows how crucial recurrent architectures are in capturing temporal dependencies in financial time series data.

Integrating semantic features along with LSTM can improve the performance of classification, which shows that sentiment provides complementary information of the market. TradeLens achieves the best performance in every metric evaluated, showing how combination of representations sequential price with sentiment features is crucial. 

Figure~\ref{fig:model_comparison} shows a comparative analysis of model performance across the evaluated metrics.

\begin{figure}[!t]
\centering
\includegraphics[width=0.5\textwidth]{pics/model_comparison.png}
\caption{Comparison of different model configurations based on evaluation metrics}
\label{fig:model_comparison}
\end{figure}

\begin{table}[!t]
\centering
\caption{Prediction samples showing correct and incorrect classifications by different models}
\label{tab:model_results}

\resizebox{\columnwidth}{!}{%
\begin{tabular}{|l|c|c|c|c|}
\hline
\textbf{Models} & \textbf{Accuracy} & \textbf{Precision} & \textbf{Recall} & \textbf{F1 Score} \\
\hline
Linear Reg. & 0.62 & 0.60 & 0.61 & 0.60 \\
Random Forest & 0.79 & 0.78 & 0.79 & 0.78 \\
LSTM & 0.84 & 0.83 & 0.85 & 0.84 \\
TradeLens & 0.87 & 0.86 & 0.87 & 0.86 \\
\hline
\end{tabular}
}
\end{table}

\subsection{Model Convergence Analysis}

The model generalizes well, as the training and validation loss has shown convergence after about 35 epochs, hence ensuring no overfitting occurs. Figure~\ref{fig:learning_curve} shows LSTM model's learning curve through 50 epochs, ensuring reduction is stable in Mean Squared Error (MSE) loss for training as well as validation sets.

\begin{figure}[htbp]
\centering

\subfloat[LSTM training and validation loss convergence over 50 epochs]{%
\includegraphics[width=0.45\textwidth]{pics/learning_curve.png}
\label{fig:learning_curve}
}
\hspace{5pt}
\subfloat[Confusion Matrix of the proposed model on the test data]{%
\includegraphics[width=0.45\textwidth]{pics/confusion_matrix.png}
\label{fig:confusion_matrix}
}

\caption{Model performance visualization: learning behavior and classification results}
\label{fig:combined_results}

\end{figure}

The behavior of classification is further checked, using the confusion matrices, as shown in Figure~\ref{fig:confusion_matrix}. Baseline models noticeably have more number of false positives and negatives, especially in situations of high volatility in market and ambiguous trend boundaries. TradeLens is able to ensure the output has a very balanced confusion matrix, with high true positives and negatives, hence is reliable and stable in real world predictions in markets.

\subsection{Ablation Study}

This ablation study focuses on evaluation of the contribution of each component using the results from Table~\ref{tab:model_results} and Figure~\ref{fig:model_comparison}.

Random Forest serves as the baseline model. It is noticeable that LSTM performs better consistently which implies the positive effect sequential learning can have in order to capture differentiating representations of price data from the history.

Using sentiment features with LSTM can give a significant boost in classification performance. This implies that alternative data sources provide excessive global information about the sentiments of the market which generally cannot be captured just by looking at price features solely.

Upon combining sentiment features with LSTM there is high gains in performance, as in case of TradeLens architecture.

Overall, this study confirms that each component is crucial in adding value to the final performance. 

%========================
\section{Discussion and Recommendations}
%========================

Based on our evaluation, the following recommendations are derived. Table~\ref{tab:tradelens_config} summarizes the recommended TradeLens configurations across different trading tiers.

\begin{table}[!t]
\centering
\caption{Recommended TradeLens Configurations for Different Market Segments}
\label{tab:tradelens_config}

\resizebox{\columnwidth}{!}{%
\begin{tabular}{lcccc}
\toprule
\textbf{Trading Tier} & \textbf{Target Scenario} & \textbf{Numerical Model} & \textbf{Explainability} & \textbf{Total Latency} \\
\midrule
Intraday & High-Frequency Signals & RF-Opt & Parallel Vector & 0.20 ms \\
Swing & Daily/Weekly & Hybrid Ensemble & RAG-Mini & 0.69 ms \\
Institutional & Portfolio Risk & Deep LSTM-120 & Full RAG & 1.25 ms \\
\bottomrule
\end{tabular}
}
\end{table}

During periods of high market volatility, combining temporal and structural features enables more reliable risk classification by leveraging LSTM and Random Forest for predictions and risk assessments.

Hybrid Ensemble Synergy model combines sequential price analysis using LSTM along with volatility and risk classification using Random Forest. The above combination benefits the model in not just capturing momentum patterns but also the structural market changes.

TradeLens is designed to ensure efficiency, while leveraging a lightweight backend and a well optimized mechanism of retrieval to ensure low latency. The architecture can be scaled for several other financial markets beyond just the Indian context.

Future work will be focused on improving the adaptability of the model using reinforcement learning and dynamic optimization of strategy, along with better sentiment analysis as per the market situations.

%========================
\section{Conclusion}
%========================

In this paper, we proposed \textit{TradeLens}. It is an artificial intelligence framework designed to improve the Indian stock trading ecosystem by closing the gap between predictive performance and institutional-grade transparency. Experimental evaluation shows that the Hybrid Ensemble Synergy Model, which combines Long Short-Term Memory (LSTM) networks with the optimized Random Forest (RF-Opt) module, was able to achieve a risk classification accuracy of 94.3\%. The system satisfies stringent institutional requirements by providing ultra-low inference latency of 0.0051~ms, which ensures lag-free and real-time decision support. TradeLens also ensures the ``black-box'' challenge is taken into consideration with deep learning models through its Retrieval-Augmented Generation (RAG) architecture.

%========================
\bibliographystyle{tfq}
\bibliography{references}

\end{document}
